\chapter{Limitaciones y amenazas a la validez}
\label{chap:limitaciones}

Este capítulo no es una disculpa genérica. Cada limitación se enuncia junto a la cifra concreta que
permite dimensionarla, porque una limitación que no puede cuantificarse no orienta al lector sobre
cuánto debe descontar de las conclusiones. La Tabla~\ref{tab:limitaciones} las reúne ordenadas por
severidad.

\begin{table}[H]
  \centering
  \footnotesize
  \caption{Limitaciones, evidencia que las cuantifica y severidad asignada.}
  \label{tab:limitaciones}
  % Tabla compuesta: cada fila enlaza una limitación con la cifra concreta que la
% cuantifica. Las cifras proceden de los artefactos citados en el capítulo:
% lo predictivo del Model Study 3 y lo económico del ganador del Portfolio Study.
\begin{tabularx}{\textwidth}{lXl}
\toprule
Limitación & Evidencia que la cuantifica & Severidad \\
\midrule
Multiplicidad de configuraciones &
Deflated Sharpe 0,682 sobre tres pasadas encadenadas más 1.728 carteras, por debajo del umbral
0,95 &
Alta \\
\addlinespace
Solapamiento de cohortes &
117 cohortes mensuales equivalen a $\approx 9$ observaciones independientes con horizonte de 12
meses &
Alta \\
\addlinespace
Potencia de la era reservada &
6 cohortes cerradas y $\approx 1$ observación independiente efectiva; 1,41 años de cartera y 1 de
2 años por encima del índice &
Alta \\
\addlinespace
Cartera elegida sobre la ventana de selección &
La ganadora es la mejor de 1.728: su IR de 0,844 en 2015--2024 es una cota superior optimista, no
una estimación insesgada &
Alta \\
\addlinespace
Dominio de un solo agente &
\textit{risk} aislado 0,1227 frente a 0,1090 del meta, que además le asigna más del 95\,\% del
peso; la neutralización retiene 86,62\,\% &
Media \\
\addlinespace
Mecanismo de la señal sin explicar &
El agente dominante se apoya en \texttt{gap\_21d} y \texttt{range\_63d} (78\,\% de las
observaciones), microestructura a semanas que ordena retornos a doce meses &
Media \\
\addlinespace
Alfa factorial no significativo en selección &
Alfa 0,27\,\% por periodo, $t=1{,}44$, $R^2=0{,}017$ &
Media \\
\addlinespace
Costes constantes &
5 pb de comisión y 10 pb de \textit{slippage}, sin impacto de mercado ni capacidad, sobre un
turnover del 324\,\% &
Media \\
\addlinespace
Universo restringido &
S\&P 500 estadounidense; 278 \textit{tickers} en 2003 frente a más de 500 en años recientes &
Media \\
\addlinespace
Versiones de catálogo heterogéneas &
Pasadas 1 y 2 bajo catálogo v6, pasada 3 bajo v7, que invierte un desempate por simplicidad &
Baja \\
\addlinespace
Salvaguarda \textit{a priori} en la curva de alfa &
No se activa en ninguna fila: la calibración usa horizonte (80\,\%) o era (20\,\%) &
Baja \\
\addlinespace
Factores no construibles &
Tamaño e inversión declarados no replicables en la regresión factorial &
Baja \\
\bottomrule
\end{tabularx}

  \caption*{La severidad es un juicio del autor sobre cuánto condiciona cada limitación las
  conclusiones, no una medida estadística.}
\end{table}

Cada una de esas limitaciones aparece ya, con su cifra, junto al resultado que acota en los
Capítulos~\ref{chap:resultados} y~\ref{chap:resultados-economicos}; la tabla las reúne para poder
verlas juntas y ordenadas. Quedan dos precisiones que ninguna fila puede expresar por sí sola.

La primera es que la multiplicidad se agravó \emph{deliberadamente}, dos veces: al encadenar tres
Model Studies y al recorrer un cartesiano de 1.728 carteras. No es un descuido sino una decisión que
asume su precio, y ese precio está declarado. Pero sus dos alcances son distintos y conviene no
confundirlos: la multiplicidad predictiva contamina la métrica que el trabajo reporta como resultado
principal, mientras que la de la rejilla de cartera no puede tocar el Rank-IC y se limita a
convertir el Information Ratio de la ventana de selección en una cota superior optimista. Ninguna de
las dos anula el contraste de permutación, que es independiente del número de configuraciones
ensayadas.

La segunda afecta al predominio de \textit{risk}, y la atribución local de la
Sección~\ref{sec:atribucion-local} lo matiza en un sentido y lo agrava en otro. Lo matiza porque ese
agente no resulta ser una réplica de la prima de baja volatilidad y porque reasigna atención entre
sus variables según el régimen, de modo que no es una apuesta congelada. Lo agrava porque desplaza
la incógnita sin resolverla: el trabajo no explica \emph{por qué} unas variables de estructura de
negociación a semanas ordenan retornos a doce meses, y una señal cuyo mecanismo no se comprende es
más frágil ante un cambio de régimen que una respaldada por una hipótesis económica contrastada.

Todo ello conduce a una lectura disciplinada: la investigación demuestra una ordenación predictiva
robusta bajo un protocolo concreto y sugiere utilidad económica, pero requiere más cohortes
cerradas, universos adicionales y un catálogo más estrecho antes de extrapolarla.

\section{Amenazas a la validez interna}

El diseño point-in-time reduce de forma directa el \textit{lookahead bias}, pero depende de la
calidad de las fechas de filing y de la correspondencia entre métricas de Finnhub y periodos
contables. Los restatements históricos pueden ser invisibles en años tempranos: que una cifra esté
hoy disponible con una fecha de filing antigua no garantiza que sea idéntica a la cifra que conoció
el mercado ese día. La mitigación consiste en conservar periodo, filing y lag de ejecución, y en no
presentar la fuente como una cinta histórica perfecta.

La pertenencia histórica al índice limita el sesgo de supervivencia de composición, pero no elimina
la cobertura desigual de empresas retiradas ni la posible pérdida de fundamentales antiguos. El
estudio mide la fracción de filas utilizables y protege contra tickers reciclados; una réplica con
una fuente comercial point-in-time sería una prueba adicional de robustez. Esta limitación afecta
principalmente a la generalización, no a la trazabilidad del panel que sí se evaluó.

\section{Amenazas económicas}

Los costes son constantes, 5 pb de comisión y 10 pb de \textit{slippage}. Esto demuestra que el
resultado no requiere coste cero, pero no modela impacto de mercado, \textit{bid--ask} dinámico,
capacidad ni fiscalidad. Dado el turnover del 324\,\% anualizado, una especificación dependiente de
liquidez podría modificar apreciablemente la cartera sin alterar el Rank-IC, y la cartera adoptada
--- ocho posiciones, sin efectivo --- es más sensible a esa especificación que la de partida.

El universo estadounidense de gran capitalización limita además la extrapolación a pequeñas
empresas, otros países o regímenes de liquidez diferentes.

La lectura correcta es específica, y conviene enunciarla porque cada limitación acota una afirmación
distinta: la multiplicidad limita el Sharpe, el solapamiento limita la precisión, la cobertura limita
la generalización y la cartera limita la transferencia de señal. Ninguna de estas salvedades permite
ignorar los placebos, la permutación o la estabilidad por eras; cada conclusión debe mantenerse
dentro del alcance de la evidencia que realmente la respalda.

\section{Limitaciones heredadas del proceso de desarrollo}

Tres limitaciones no proceden del diseño sino de decisiones tomadas durante el desarrollo. Merecen
figurar aquí porque afectan a la interpretación de las cifras y no serían visibles para un lector
que sólo viera el resultado final; su origen y la disyuntiva técnica que las motivó se documentan en
el Anexo~\ref{anexo:desarrollo}.

La primera es la salvaguarda de la curva de alfa: una recta fija entre $-10$\,\% y $+10$\,\% que se
aplicaría si ninguna ventana de la cascada produjese una pendiente creciente. En el estudio de
referencia \textbf{no se activa en ninguna fila}: la calibración usa la ventana de horizonte en el
80\,\% de los casos y la de era en el 20\,\% restante. Deja de ser, por tanto, una limitación
material de estas cifras y pasa a ser una decisión de diseño discutible que conviene declarar ---
un supuesto \emph{a priori} que se dispararía precisamente cuando la evidencia indica que no hay
señal ---, pero que ninguna cifra de este trabajo utiliza.

La segunda es que la calibración isotónica que precedió a esa salvaguarda se retiró sin conservarla
como opción del catálogo, de modo que este trabajo \emph{no puede} ofrecer una comparación empírica
entre ambas formulaciones; la afirmación de que la nueva es mejor descansa en un argumento sobre el
estimador, no en una medición pareada.

La tercera afectaba a las variables de cartera, que en versiones anteriores de este trabajo tomaban
su valor del catálogo recomendado y no de una medición, de modo que las cifras económicas debían
leerse como una cota inferior. El Portfolio Study resuelve esa limitación y la sustituye por otra de
naturaleza distinta: la cartera ya no es arbitraria, pero es la mejor de 1.728 evaluadas sobre la
ventana de selección, y por tanto sus cifras dentro de esa ventana son ahora una cota \emph{superior}
optimista en lugar de una inferior conservadora. La comprobación de la era reservada es la única
lectura de esas cifras que no está sujeta a ese sesgo, y descansa sobre poco más de un año de datos.

\section{Reproducibilidad y versiones de catálogo}

Las tres pasadas de la cadena no corrieron bajo la misma versión del catálogo, lo que las hace no
estrictamente comparables en sus reglas de desempate; el Anexo~\ref{anexo:catalogo} detalla qué
cambió y por qué no altera ninguna medición de Rank-IC.

Lo que sí condiciona la lectura de los resultados es la consecuencia: dos de las diecisiete
decisiones de la pasada de referencia se resolvieron por empate técnico, es decir, por una
preferencia declarada de antemano y no por medición. Sus valores finales
--- \artefacto{execution\_lag\_days} y \artefacto{lgbm\_min\_child\_samples} --- no deben
presentarse como hallazgos empíricos. En el segundo caso, además, la alternativa medía
\emph{ligeramente mejor} que el valor elegido.
